\documentclass[11pt]{article}
\usepackage[margin=1in]{geometry}
\usepackage{booktabs}
\usepackage{array}
\usepackage{xcolor}
\usepackage{hyperref}
\usepackage{enumitem}
\hypersetup{colorlinks=true,linkcolor=blue!50!black,urlcolor=blue!50!black}

\title{\textbf{HyMeKo Exploration Summary}\\[3pt]
\large Results, documents, research venues, and reachable potential}
\author{Codex summary for Dr. Csaba Hajdu}
\date{2026-06-30}

\newcommand{\pathref}[1]{\nolinkurl{#1}}
\newcommand{\good}{\textcolor{green!45!black}{strong}}
\newcommand{\mixed}{\textcolor{orange!70!black}{mixed}}
\newcommand{\open}{\textcolor{blue!60!black}{open}}

\begin{document}
\maketitle

\begin{abstract}
This document consolidates the broad repository exploration performed on
2026-06-30. The short conclusion is that HyMeKo is now best understood as three
things at once: a mature signed-hypergraph infrastructure stack, a disciplined
experimental lab for testing whether structure helps learning and control, and
a paper/evidence production system. The cleanest near-term contribution is the
machine-verified cross-view consistency kernel. The highest-upside research
lines are topology/arity matching for control, multiplicative holonomy/transport,
and the new Fiber-Spike-Rotor language model. The RL line is promising but not
yet a clean ``HSiKAN wins control'' story.
\end{abstract}

\section{Executive read}

\textbf{The strongest current claim} is not that HSiKAN beats every baseline.
The strongest claim is that HyMeKo provides one canonical signed-hypergraph
source of truth, emits many concrete views, and machine-checks that invariants
recovered from those emitted views agree.

That claim is backed by the cross-view consistency work:
\begin{itemize}[nosep]
  \item \pathref{reports/2026-06-28-cross-view-consistency.md}
  \item \pathref{reports/2026-06-29-acceptance-novelty-execution.md}
  \item \pathref{verification/cross_view_consistency/}
\end{itemize}

The research layer is more volatile but unusually healthy: negative results are
recorded, confounds are named, and the best follow-ups are discriminator tests
rather than ``train larger'' guesses.

\section{What the repository contains}

\begin{table}[h]\centering\small
\begin{tabular}{p{0.25\linewidth}p{0.65\linewidth}}
\toprule
Area & Observed state \\
\midrule
Rust framework & Workspace crates for parser, core IR, query engine, transforms, HRE,
emitters, WASM/editor, server, MCP, wire protocol, formats, monitoring, graph,
P-graph, Clifford, and Python bindings. \\
Evidence system & Plans under \pathref{docs/plans/}; acceptance records under
\pathref{reports/}; milestone indexes in \pathref{DONE.md}, \pathref{BACKLOG.md},
and \pathref{ROADMAP.md}. \\
Transforms and views & URDF, SDF, MJCF, DOT, Mermaid, ROS2 launch, SysML-style
requirements views, Gazebo/MuJoCo-related emitters, torch-dataflow templates. \\
Learning side & Signed-link, HSiKAN/SignedKAN/Gomb, Soma/RicciStim vision,
HyMeYOLO, RL/control, StructuralActor, and the new \pathref{hymeko_lm/}. \\
Artifacts & Benchmark CSVs, checkpoints, generated figures, GIFs, PDFs, LaTeX
reports, and experiment JSON/JSONL logs. \\
\bottomrule
\end{tabular}
\end{table}

\section{Mature infrastructure}

The framework layer looks mature enough to support papers and demos. The
architecture is documented in \pathref{architecture/README.md}; the workspace
uses a parser/IR/query/emit separation; benchmark CSVs in
\pathref{hymeko_bench/results/} show compilation and expansion measurements;
and the cross-view harness drives the real CLI emitters rather than a toy
emitter.

The most valuable engineering property is the hub-and-spoke model:
\[
  X_f(\epsilon_f(H)) = Q(H)
\]
where an emitted concrete view is parsed back into a common invariant and
checked against the IR-level query. This turns ``views should agree'' into a
testable commuting-square claim.

\section{Core results}

\begin{table}[h]\centering\small
\begin{tabular}{p{0.22\linewidth}p{0.50\linewidth}p{0.18\linewidth}}
\toprule
Line & Current finding & Status \\
\midrule
Cross-view IR & 16 robotics fixtures across five views, plus systems-engineering
requirements fixtures; drift-prevention demo; Z3/sympy proof support. & \good \\
Topology matching & Matched controller topology wins 8/8 supervised plant
targets; dense complete graph is not a universal winner. & \good \\
Arity matching & Hypergraph lift hurts graph plants but wins decisively on
hypergraph plants. Match controller arity to plant arity. & \good \\
Signed-link & Strict/leakage-audited line is valuable; rotor geometry itself was
not load-bearing in AUROC ablations. Structural features and protocol discipline
are the real contribution. & \mixed \\
Vision/Soma & Mean-pool falsification was partly a readout artifact. Spatial
pyramid/tree readout gives a scalable position-aware compromise. & \mixed \\
FSR language model & TinyShakespeare PoC: FSR pre-norm beats matched small
transformer on bpb; sparse top-k fixes the quadratic wall; sphere residual loses. & \open \\
RL/control & Infrastructure is serious, BC works, PPO success collapses, and no
clean HSiKAN-over-MLP control win is established yet. & \open \\
\bottomrule
\end{tabular}
\end{table}

\section{Vision line}

The vision story is no longer a simple failure or success. Earlier mean-pool
Soma/RicciStim results underperformed conventional controls. The June 29--30
readout work changes the diagnosis: the bottleneck was often positional
compression.

Key result from \pathref{reports/2026-06-29-soma-spatial-tree-readout.md}:
\begin{center}\small
\begin{tabular}{lrr}
\toprule
Readout & Parameters & Cluttered-MNIST accuracy \\
\midrule
Mean-pool & 2010 & 0.314 \\
Attention & 2027 & 0.466 \\
Spatial tree & 5227 & 0.537 \\
Flatten & 24890 & 0.617 \\
\bottomrule
\end{tabular}
\end{center}

The useful principle is that a flatten readout over a spatial field can be
compressed into a multi-scale spatial-pyramid pool: roughly one fifth the
parameters for roughly ninety percent of the accuracy, with scale and
variable-anchor support. This is exportable to HyMeYOLO heads, RL vision, and
positioned graph/point readouts.

\section{FSR language model}

The new \pathref{hymeko_lm/} package is a clean-slate language-model experiment.
The Fiber-Spike-Rotor mixer replaces softmax attention with a causal,
signed, rotor-transported, spike-gated gather. The channel mixer uses HSiKAN/CR
or Chebyshev-CR style machinery.

Reported Phase 1 result in \pathref{reports/2026-06-29-fsr-lm-phase1.md}:
\begin{center}\small
\begin{tabular}{lrr}
\toprule
Model & Validation bpb & Throughput \\
\midrule
Transformer control & 2.570 & 35.1k tok/s \\
FSR sphere & 2.770 & 11.8k tok/s \\
FSR pre-norm & 2.479 & 11.2k tok/s \\
\bottomrule
\end{tabular}
\end{center}

The honest interpretation is important. The positive result belongs to the FSR
mixer plus HSiKAN channel mix with a standard pre-norm residual. The Gomb sphere
residual does not currently earn its place on this task. The speed report
\pathref{reports/2026-06-29-fsr-lm-speed.md} adds hard top-k sparsity:
\pathref{spike_k=16} preserves or improves quality and makes long context
feasible, but absolute throughput still trails fused transformer attention.

\section{Holonomy}

The holonomy-discriminator line is the sharpest recent scientific update. The
initial plan was to test whether cycle holonomy is load-bearing. The current
implementation in \pathref{hymeko_rl/holonomy_probe.py} distinguishes additive
signed-walk aggregation from multiplicative transport.

Verification run performed during this summary session:
\begin{center}\small
\begin{tabular}{lrr}
\toprule
Arm & Accuracy & Parameters \\
\midrule
Transport / rotor-like multiplicative reader & 1.000 & 49 \\
Additive signed-walk reader & 0.4805 & 225 \\
MLP & 0.5913 & 1505 \\
Linear confound guard & 0.5054 & 13 \\
\bottomrule
\end{tabular}
\end{center}

This is a useful but sobering result. It supports the idea that pure cycle
holonomy is multiplicative/transport-like, but it downgrades expectations that
plain additive HSiKAN message passing will solve parity-style holonomy tasks.

\section{RL and control}

The RL part is promising, but it is not yet a clean ``HSiKAN wins RL'' story.
It should be described as an experimental substrate plus a developing theory of
when structure should matter.

\subsection*{What is built}

\begin{itemize}[nosep]
  \item MuJoCo/Gym-style tasks: reaching, pick-and-place, Galambos/collaborative,
  quadruped/walker variants, sanity worlds.
  \item PPO, BC warm-starts, off-policy experiments, rendering, checkpoints, and
  policy builders for MLP, HSiKAN, SA-HSiKAN, and mixtures.
  \item Fast testbeds in \pathref{hymeko_rl/sanity_rl.py} and
  \pathref{hymeko_rl/sanity_worlds.py}: bandit, grid, hex, and collaborative.
  \item A pre-screen entry in \pathref{hymeko_rl/rl_prescreen.py}.
\end{itemize}

\subsection*{Current empirical state}

From \pathref{reports/2026-06-29-rl-overnight.md}, BC is the reliable worker:
pick-and-place after BC reaches useful lift/place success. PPO then improves
scalar return while collapsing actual task success to near zero. That is a
reward-alignment problem, not yet an architecture win.

\begin{center}\small
\begin{tabular}{lcc}
\toprule
Policy & After BC lift/place & After PPO lift/place \\
\midrule
HSiKAN & 0.875 / 0.625 & 0.000 / 0.000 \\
MLP & 1.000 / 0.500 & 0.125 / 0.000 \\
\bottomrule
\end{tabular}
\end{center}

The right near-term RL claim is therefore not performance superiority. It is:
\begin{quote}
HyMeKo provides a structural, declarative substrate for control experiments; the
standing hypothesis is that structure helps when the controller topology/arity
matches the plant and task invariant.
\end{quote}

\subsection*{Main blockers}

\begin{enumerate}[nosep]
  \item \textbf{PPO success collapse.} Scalar return can improve while success
  gets worse. The fix path is BC anchoring, success-aligned terminal rewards,
  and reporting task success separately from return.
  \item \textbf{Launch-bound HSiKAN.} Vanilla HSiKAN is slow in rollout because
  the forward is many tiny tensor operations. SA-HSiKAN and structural collapse
  are the right directions.
  \item \textbf{Wrong tasks can hide structure.} Linear/local tasks let MLP tie.
  Holonomy, contact cycles, gait, and matched topology are better discriminators.
\end{enumerate}

\subsection*{Best RL moves}

\begin{itemize}[nosep]
  \item Use the fast toy lab before every slow MuJoCo run.
  \item Lead demos with BC policies until PPO refinement stops degrading success.
  \item Use SA-HSiKAN for rollout-heavy tasks.
  \item Validate topology matching in closed loop, not only supervised structural
  proxies.
  \item Reserve rotor/transport backbones for cyclic or holonomy-heavy tasks.
\end{itemize}

\section{Best venues}

\begin{table}[h]\centering\small
\begin{tabular}{p{0.25\linewidth}p{0.65\linewidth}}
\toprule
Venue family & Best matching story \\
\midrule
IEEE T-SMC / IEEE SMC & Canonical hypergraph IR, machine-verified cross-view
consistency, robotics plus SysML witnesses, drift-prevention demo. \\
Formal methods / MDE & Consistency by convention versus machine-verified
commuting diagram over real emitters. \\
Control / CPS & Topology and arity matching: the controller should match the
plant's structural coupling, not merely have more edges or parameters. \\
Graph ML / signed networks & Leakage-audited structural features, strict
protocols, and inductive signed-link learning. \\
ML architecture workshop & FSR language model with rotor transport, spike top-k,
and HSiKAN channel mix; must scale beyond TinyShakespeare before strong claims. \\
Vision / neuro / perception & Spatial-pyramid readout, RicciStim structure
diagnostics, and possible brain-predictivity rather than raw accuracy claims. \\
Robotics collaboration & Declarative control scenarios, BC-anchor refinement,
matched topology, affect/reward models, and Kato-style collaborative grasping. \\
\bottomrule
\end{tabular}
\end{table}

\section{Recommended next sequence}

\begin{enumerate}
  \item \textbf{Ship the cross-view consistency paper first.} It is the cleanest
  evidence package and the strongest novelty claim.
  \item \textbf{Turn topology/arity matching into a closed-loop control result.}
  This is the natural next high-value RL/control discriminator.
  \item \textbf{Make multiplicative holonomy the gauge mechanism.} Do not expect
  additive HSiKAN alone to solve parity-like holonomy.
  \item \textbf{Repair RL refinement.} BC anchor plus success-aligned reward must
  precede more architecture claims.
  \item \textbf{Scale FSR carefully.} Lead with \pathref{spike_k=16}; treat Gomb
  sphere as currently negative on byte-LM.
  \item \textbf{Extract and reuse SpatialPyramidPool.} It is a practical component
  with obvious export routes.
\end{enumerate}

\section{Verification performed for this summary}

The first attempt used system Python and failed because \pathref{torch} and
\pathref{pytest} were unavailable. Retrying through the project virtual
environment succeeded:
\begin{verbatim}
.\.venv\Scripts\python.exe -m hymeko_rl.holonomy_probe --ring-size 12 --seeds 2
.\.venv\Scripts\python.exe -m pytest -q ^
  hymeko_rl/tests/test_holonomy_probe.py hymeko_rl/tests/test_rl_prescreen.py
\end{verbatim}

Result: transport solved the holonomy probe, additive/MLP/linear stayed near
chance, and the two test files passed: \textbf{4 passed in 14.85 s}.

\section{Final judgement}

HyMeKo has one publishable formal-methods core now, several high-ceiling research
programs, and a useful habit of recording falsification. The highest potential is
not in claiming universal structural superiority. It is in showing exactly when
structure is load-bearing: verified cross-view invariants, matched controller
topology, matched arity, and multiplicative holonomy transport.

\end{document}
